\section{Scientific Activity}\label{scientific-activity}

My research activity has been developed in the broad area of robotics
perception and control, with applications to relevant societal problems.
I study the fundamental problem of how intelligent systems (natural and
artificial) can make \textbf{use of its sensors, actuators, and
computing resources to learn and adapt to unstructured and dynamic
environments}. Then, I try to apply my findings to improve current
robotics technologies and have a better understanding of natural
intelligent systems. My work has been applied in several areas of
societal impact such as \textbf{Health \& Care}, \textbf{Environmental
Monitoring} (Wildfire and Maritime scenarios), \textbf{Humanoid
Robotics}, \textbf{Security and Surveillance}.

My research spans multiple disciplines at the intersection of
\textbf{computer vision, robotics, machine learning, cognitive systems,
optimization, and control}. Most of my research lines are carried-out in
collaboration with colleagues of the \textbf{Computer and Robot Vision
Laboratory} (Prof. Santos-Victor, Prof. José Gaspar, Prof.
Salvador-Marques, Prof. Catarina Barata, Prof. Bruno Damas, Prof. Nuno
Pessanha-Santos, Dr. Plinio Moreno). Additionally, I have created links
with many researchers worldwide, both on the engineering and the natural
sciences. As a member of ISR-Lisboa, I have joined the Associated
Laboratory \textbf{LarSyS (Laboratory of Robotics and Engineering
Systems)}, where I coordinate the Thematic Line ``\textbf{INTERACTION -
Cognitive Robots and Human Experience}''.

\subsection{Research Lines}\label{research-lines}

In the following paragraphs I describe in more detail in my core
research activities. I will refer to related research projects, PhD
students, and relevant collaborations. The actual PhD supervision teams
and project's details are described in Appendices \ref{supervision-activities} and \ref{project-grants}.

\subsubsection{Bioinspired and Learning
Robots}\label{bioinspired-and-learning-robots}


One of my main research lines is on ``Bioinspired and Learning Robots''.
I have completed my PhD in 2004, on the topic of ``Binocular Head
Control with Foveal Vision'', with inspiration on the mechanics and
physiology of the human visual system. In 2006, I joined the
\textbf{ROBOTCUB Project}, whose purpose was the development of a
standard robotic platform for research on developmental and cognitive
robotics and involved a multidisciplinary research team from the
engineering sciences and the human sciences (neuroscience and
psychophysics). In that pioneering project, I have developed researched
work on humanoid technologies, visual attention, active vision, object
tracking, and learning object affordances. I had the opportunity to
participate in the joint effort of the development of \textbf{the iCub
Robot}, one of the most complex humanoid robots worldwide and a unique
platform in Portugal. A fascinating topic on those years, prior to deep
learning technology, was object recognition. In my first PhD student
co-supervision (Dr. Plinio Moreno, with Prof. José Santos-Victor), we
researched on biologically inspired computer vision for object
recognition. On the follow-up of the ROBOTCUB project I participated in
projects \textbf{MIRROR, CONTACT and ROBOSOM}, also focused humanoid
technologies through synergies from Robotics and Neurosciences. Also, I
have been awarded two national projects, \textbf{BIOLOOK} (2007-2010)
and \textbf{BIOMORPH} (2014-2015). In those projects I have researched
eye-head coordination algorithms and visual prediction mechanisms for
motion estimation (PhD of Dr. Nuno Moutinho, with Prof. José Gaspar) and
studied the development of an artificial fovea through the optimization
of sensorimotor learning principles (PhD of Dr. Jonas Ruesch, with Prof.
José Santos-Victor). Integrated in the IST-EPFL doctoral program on
\textbf{Robotics, Brain and Cognition} (RBCog), I engaged in
collaborations with Prof. Auke Ijspeert from the Biorobotics Laboratory
and EPFL, Lausanne, Switzerland, in bioinspired locomotion behaviors in
quadrupeds, and modular robots, that led to the graduation of 2 PhD
students (Dr. Alexandre Tuleu, Dr. Mehmet Mutlu). Also in the RBCog
Doctoral Program, I collaborated with Prof. Helder Araújo (Univ.
Coimbra) on space-variant vision mechanisms for resource constrained
systems (Ph.D of Dr. Rui Figueiredo), and collaborated with Prof.
Lorenzo Jamone (QMUL) and Prof. José Santos-Victor on humanoid robot
self-supervised learning of body schema (PhD of Dr. Pedro Vicente) which
was awarded the prize of best 2021 national PhD thesis in Robotics by
the Portuguese Society of Robotics. More recently, I have been
collaborating in the Horizon 2020 project \textbf{ORIENT}, an ERC
initiative of Prof. John van Opstal from the Donders Institute for
Brain, Cognition and Behaviour, Radboud University, Nijmegen,
Netherlands. In this project the ISR/IST team developed a biomimetic
robot eye to study the principles underlying the optimization of the
fast ocular movements called saccades.

\subsubsection{Vision for Decision and
Control}\label{vision-for-decision-and-control}


Another of my main research lines is on computer vision technologies for
decision and control, such as motion estimation, activity recognition,
object detection, classification, and tracking. While still working on
my PhD, I joined the \textbf{NARVAL} Project, where I assumed an active
role in the development of visual navigation behaviors for underwater
autonomous systems. These are fundamental skills for the autonomy of
advanced robotics and surveillance systems. In the sequence, I worked on
probabilistic methods for motion analysis for autonomous vehicles (PhD
of Dr. Hugo Silva) in collaboration with Prof. Eduardo Silva (ISEP-IPP).
I established collaborations with the Video Surveillance Industry
through the innovation projects \textbf{HDA} (ANI) and \textbf{DICORE2S}
(EU-FP7) that led to novel methods in the detection of pedestrians (Dr.
Matteo Taiana, with Prof. Jacinto Nascimento), multiple-hypothesis
tracking in camera networks (Dr. Dario Figueira, with Prof. Jacinto
Nascimento), and gait-based person re-identification (Dr. Athira
Nambiar, with Profs. Jacinto Nascimento and José Santos-Victor). In the
context of the CMU-Portugal Program, I have co-supervised Ph.D. students
in the development of novel image classification algorithms using matrix
completion methods (PhD thesis of Dr. Ricardo Cabral, IBM Prize 2014) in
collaboration with Prof. João Paulo Costeira (ISR/IST) and Prof.
Fernando de La Torre (CMU), and on activity recognition from videos with
Deep Learning methods (PhD Thesis of João Antunes Martins) in
collaboration with Prof. Dan Siewiorek (CMU) and Prof. Asim Smailagic
(CMU). I have made collaborations with the Portuguese Navy and Air Force
through national research projects on maritime surveillance
(\textbf{SEAGULL}, 2013-2015), airborne image analysis of maritime and
forest environments (\textbf{VOAMAIS}, 2019-2023), and wildfire
monitoring (\textbf{FIREFRONT}, 2019-2023) where I supervised PhD
students in 3D model-based tracking of Drones (thesis of Dr.
Pessanha-Santos) in collaboration with Prof. Vitor Lobo (Navy Academy),
and 2D tracking of vessels in maritime environments (Dr. Gonçalo Cruz)
in collaboration with Prof. António Pedro Aguiar (FEUP). More recently I
have established collaborations with Prof. Alessandro Saccon (TUe,
Eindhoven), through the co-supervision of MSc students, to further
research 3D model-based tracking algorithms for robotic manipulation
tasks.

\subsubsection{Object Grasping and
Manipulation}\label{object-grasping-and-manipulation}

Object grasping and manipulation are fundamental robotic skills for
industry and are quite complex to deploy in unstructured environment.
Object grasping and manipulation rely on physical contacts with objects
that involve impacts with multiple surfaces in simultaneity, being
extremely complex to model and simulate, particularly with human-like
robot hands. Object grasping and manipulation relate to computer vision
as the initial phases of grasping can be guided by vision, but as soon
as fingers contact the objects, tactile and force sensing become
essential and, afterwards, state estimation and control can fuse vision,
force, and tactile sensing. After my participations in the
\textbf{ROBOTCUP} and \textbf{MIRROR} Projects, where I researched
visual skills involved in robotic manipulation, I participated actively
in the resubmission of a project proposal for research in dexterous
robotic manipulation (\textbf{HANDLE} project), in which I assumed the
role of Technical Manager. In collaboration with the Principal
Investigator, Prof. Véronique Perderaux from the ISIR lab of UPMC,
Paris, I researched human dexterity and developed methods for the
high-level control of in-hand manipulation (PhD of Dr. Urbain Prieur).
In comparison to visual sensors, which are becoming close to human
resolution, tactile sensors for robot hands are still many orders of
magnitude behind.

In the context of the individual Marie Curie grant of Dr. Lorenzo
Jamone, the \textbf{LIMOMAN} project developed magnetic tactile sensors
for the instrumentation of dexterous robotic hands and several methods
were developed for tactile based control of robotic grasps. To further
develop the tactile sensing technology, we established collaborations
with Prof. Susana Freitas (IST, INESC-MN) and co-supervised two PhD
students (Dr. Pedro Ribeiro and Dr. Miguel Neto) that developed highly
sensitive magneto-resistive sensors and applications to robotics.

\subsubsection{Human-Robot Interaction}\label{human-robot-interaction}

My research in Human-Robot Interaction attempts to deploy in robots many
of the technologies developed in my other research lines, and evaluate
how such systems can be used, valued, and enjoyed by humans. In
traditional robotics applications, robots perform their tasks
autonomously and independently from humans, most often in separate
physical spaces due to safety constraints. As robots become more
versatile, robust, and safe, they can be used in applications involving
close interaction with humans, e.g., domestic use, collaborative
manufacturing, and services. However, to be accepted by humans, beyond
competence and reliability in executing the required tasks, such robots
must provide a range of other features that are much related to human
psychology: usability, pleasant appearances and motions, natural
communication and behaviors, sense of safety, shared agency, and
control, among others. In my research, I have addressed many of these
dimensions. In the Poeticon++ project I have studied how object's
affordances and goal directed actions can be used by robots to
understand human intentions and provide more natural non-verbal
communication channels (PhD of Giovanni Saponaro). In the \textbf{AHA}
and \textbf{IntelligentCare} projects I have studied how interactive
robots and virtual/augmented reality systems can be used in conjunction
with therapists to aid in physical and cognitive rehabilitation
processes (PhD of Min Hun Lee, PhD of Ana Rita Cóias, ongoing). In the
PhD work of João Avelino, in co-supervision of Prof. Leonel
Garcia-Marques (Fac. Psychology, UL) and Prof. Rodrigo Ventura (ISR/IST)
we have implemented human-like salutation and engagement behaviors and
illustrated their validity in human-robot interactions ``in the wild''
(unconstrained environments). In \textbf{ShiftHRI} project, in
collaboration with Prof. Tiago Guerreiro (FCUL) and Prof. Jodi Forlizzi
(CMU) we are studying how telepresence and social robots can be used to
provide valuable services to older people (PhD of Hugo Simão, ongoing).
In the \textbf{HAVATAR} project we are studying how telepresence and
service robots can be used in a hospital context, providing value to
customers and patients' relatives.

I was the principal investigator of projects funded by the Portuguese
Science Foundation, by the Portuguese Innovation Agency, and by special
programs of collaboration of Portuguese and International Universities
(CMU-PT). Also, I was technical manager and local leader of projects of
programmes EU-FP7 and HORIZON 2020.

\subsection{Creation And Reinforcement of Research
Infrastructures}\label{creation-and-reinforcement-of-research-infrastructures}

I am the deputy coordinator of the Computer and Robot Vision Lab
(Vislab) of the Institute for Systems and Robotics (ISR), partly
responsible for the acquisition, maintenance, and organization of
general-purpose lab resources (computers, cameras, accessories,
consumables, etc). In the context of national and international research
projects I have been directly involved in the construction and
acquisition of specialized platforms and equipment that has reinforced
the Vislab infrastructure for top level research in Computer and Robot
Vision:

\begin{enumerate}
\def\labelenumi{\arabic{enumi}.}
\item
  The \textbf{iCub robot} is one of the most advanced humanoid platforms
  in the world, and unique in Portugal, as an outcome of the RobotCub
  project. This platform has been used in top-tier research projects.
\item
  The \textbf{NAO} and \textbf{Darwin-OP} humanoid robots are small
  humanoid platforms that have simple operation and programming
  interfaces to allow rapid prototyping of demonstrations and user
  studies.
\item
  The \textbf{CamNet} in a distributed camera network installed in the 3
  floors of ISR (approved by the National Data Protection Commission)
  that has been used in video surveillance research projects. A dataset
  has been acquired, that counts with more than 300 users worldwide.
  This was an outcome of the HDA project in collaboration with the
  national industry.
\item
  The \textbf{Tobii EyeTracker} is an eye-tracker device to measure the
  instantaneous gaze direction of a human subject. It was acquired in
  the context of the Biolook project and has been used for several
  psychophysics studies on human visual cognition.
\item
  The \textbf{Vizzy} robot is a unique humanoid social robot jointly
  developed through several integrative efforts. I have participated in
  the development of several hardware and software components.
\item
  The \textbf{Optitrack} high speed motion capture system allows the
  precise acquisition of trajectories and has been used to benchmark
  several robot and computer vision methods.
\item
  The \textbf{BioRob} robot in a robotic arm composed of elastic
  actuators that is capable of highly dynamic motions and was used to
  the study dynamic motion primitives and real-time trajectory control.
\item
  The \textbf{Medusa} robot head is a humanoid robot head with 4
  degrees-of-freedom where I have worked for the great part of my MSc
  and PhD and have contributed to the development of the hardware,
  control electronics, and software.
\item
  The \textbf{EuroHead} is a high precision robotic stereo head for
  metrology application, for which I have contributed for the
  specification and acquisition.
\item
  \textbf{PEPE} is an augmented reality platform for physical and
  cognitive stimulation composed of a mobile base, an ultra-wide throw
  projector, a computer and a motion sensor that was developed in the
  AHA project. It has spawned other projects and has raised interest for
  commercial exploitation. A national patent has been granted in 2024,
  and international patents submitted.
\item
  \textbf{Temi} and \textbf{SanBot} are commercial teleoperation robots
  with a high degree of autonomy that have been acquired in the context
  of the AHA and HAVATAR projects for research in robotics
  tele-autonomy.
\item
  The RBCog lab is a research infrastructure included in the National
  Roadmap of Research Infrastructures where I participate. I have
  contributed for its development through specifications. It includes
  many of the resources listed above and others like the \textbf{Baxter
  Robot}, a \textbf{Kinova Gen 3} robot arm, and manipulator
  end-effectors for several applications.
\end{enumerate}

The available infrastructure has been key in advanced training at the
level of MSc and PhD projects, as well as promoting research with
external partners that have led to long term collaborations and joint
project proposals.

\subsection{Scientific Publications}\label{scientific-publications}

I have a very diversified publication profile and publish both in
journals of my core research fields of \textbf{robotics, computer
vision, control, and artificial intelligence} (e.g. \emph{Science
Robotics, IEEE Transactions on Cognitive and Developmental Systems, IEEE
Transactions on Image Processing,} \emph{Neurocomputing, Pattern
Recognition, Neural Networks, The International Journal of Robotics
Research, Autonomous Robots, Pattern Analysis and Machine Intelligence,
Journal of Field Robotics, IEEE RA-L, IEEE Transactions on Systems, Man
and Cybernetics, Computer Vision and Image Understanding}), and in
journals of applications of robotics and AI to: \textbf{health and
rehabilitation} (\emph{International Journal of Clinical and Health
Psychology}, \emph{IEEE Transactions on Neural Systems and
Rehabilitation Engineering, Journal of Neuroengineering and
Rehabilitation}); \textbf{sensing and instrumentation} (\emph{IEEE
Internet of Things Journal, Computers and Electronics in Agriculture,
International Journal of Wildland Fire, Expert Systems with
Applications, Experiments in Fluids, IEEE Transactions on Circuits and
Systems for Video Technology, Remote Sensing, IEEE Transactions on
Geoscience and Remote Sensing, IEEE Transactions on Magnetics, Mechanism
and Machine Theory}); and \textbf{human-machine interaction}
(\emph{International Journal of Social Robotics, User Modeling and
User-Adapted Interaction, Proceedings of the ACM in Human-Computer
Interaction, ACM Computing Surveys}).

Overall, I have published \textbf{404 peer reviewed articles}
(\textbf{112 international journal papers}, \textbf{11 book chapters},
\textbf{207 papers in international conferences}, \textbf{41 papers in
international workshops and symposia}, \textbf{33 papers in national
publications}). According to Scopus, I have \textbf{463 co-authors} in
scientific papers, \textbf{6288 citations} and an \textbf{H-index of
35}, as of 26 May 2026.

The full list of my publications is provided in Appendix \ref{publications}. The impact
of my publications is analysed in Appendix \ref{impact-factors} through the publication
citations and journal impact factors provided by several indexing
services: Clarivate's Web of Science (WoS), SCImago, Elsevier's Scopus,
and Google Scholar (GS).

\subsubsection{The Ten Most Representative
Works}\label{the-ten-most-representative-works}


In the next pages I list and briefly describe the 10 works I consider most
representative about my research and that are having a significant
impact in the research community as measured by the number of citations
in WoS/GS as of May 27\textsuperscript{th} 2026. I have selected 10
works that were published in 10 different journals of the first quartile
(Q1) according to WoS Journal Citation Report (JCR) of 2022. Each of the
articles is classified into one or more of my main research domains:
Bioinspired and Learning Robots, Vision for Decision and Control, Object
Grasping and Manipulation, Human-Robot Interaction. For each article, a
couple of other related papers is included.

\clearpage

\textbf{\textsc{\ul{The iCub Robot}}}

\textbf{Authors:} G. Metta, L. Natale, F. Nori, G. Sandini, D. Vernon,
L. Fadiga, C. von Hofsten, J. Santos-Victor, \ul{A. Bernardino}, L.
Montesano

\textbf{Year:} 2010

\textbf{Title:} ``The iCub Humanoid Robot: An Open-Systems Platform for
Research''

\textbf{Journal:} \emph{Neural Networks,} 23(8-9), 1125--1134.

\textbf{Domain:} Bioinspired and Learning Robots

\textbf{IF:} 7.8

\textbf{Citations:} 469/773

\textbf{DOI:} \url{https://doi.org/10.1016/j.neunet.2010.08.010}

\textbf{Description:} This paper presents the humanoid robot iCub, in
its most relevant aspects. The iCub robot was the main output of the EU
project RobotCub. The iCub Robot is one of the most sophisticated
humanoid robots in the world. It has a child-like morphology and was
specially designed for research in developmental robotics. Currently
there are about 50 copies of this robot in the world, and one is in
Portugal (Vislab/IST/IST). The robot has been used in several PhD and
MSc theses at Vislab, and on ``spin-off'' projects like the Poeticon++
project in which Vislab participated. The iCub robot has been
fundamental in the creation of a community associated to the research of
humanoid robots, cognitive systems, and robot learning. The open-source
nature of the platform (hardware and software) promotes the exchange of
experiences among several research groups. A large part of the paper
presents results obtained by the Vislab team, that had a significant
impact in the project, like the construction of the robot head, the
design of the cognitive architecture, and the development of robot
vision and learning algorithms.

\textbf{Related Papers:}

\begin{itemize}
\item
  J. Ruesch, M. Lopes, \ul{A. Bernardino}, J. Hörnstein, J.
  Santos-Victor, R. Pfeifer (2008), ``Multimodal Saliency-Based
  Bottom-Up Attention A Framework for the Humanoid Robot iCub'',
  \emph{IEEE International Conference on Robotics and Automation (ICRA
  2008)}, pp 962-967. \url{https://doi.org/10.1109/ROBOT.2008.4543329}
\item
  R. Beira, M. Lopes, M. Praça, J. Santos-Victor, \ul{A. Bernardino}, G.
  Metta, F. Becchi, R. Saltarén (2006), ``Design of the Robot-Cub (iCub)
  Head'', \emph{IEEE International Conference on Robotics and Automation
  (ICRA)}, pp 94-100. \url{https://doi.org/10.1109/ROBOT.2006.1641167}
\end{itemize}

\clearpage

\textbf{\textsc{\ul{Object's Affordances}}}

\textbf{Authors:} L. Montesano, M. Lopes, \ul{A. Bernardino}, and J.
Santos-Victor

\textbf{Year:} 2008

\textbf{Title:} ``Learning Object Affordances: From Sensory Motor Maps
to Imitation''

\textbf{Journal:} \emph{IEEE Transactions on Robotics,} 24(1), 15-26.

\textbf{Domain:} Bioinspired and Learning Robots, Object Grasping and
Manipulation

\textbf{IF:} 7.8

\textbf{Citations:} 281/548

\textbf{DOI:} \url{https://doi.org/10.1109/TRO.2007.914848}

\textbf{Description:} According to several theories on cognition, the
way humans perceive objects is related more to their utility (what they
afford to - affordances) rather than their identity. This contrasts with
the classical view of object's representations (3D reconstruction,
recognition) in which object's models are created without taking the
task into account. Affordance is a very powerful concept for general
purpose and adaptive robots because it allows to consider object
manipulation as a pre-categorical skill: robots do not need to recognize
the semantic class of an object to manipulate it, and thus, can
experiment it for the execution of tasks even if not seen previously.
This paper shows, through a set of experiments with a humanoid robot,
that pre-categorical cues like shape features are sufficient to
manipulate the objects in simple tasks. The robot executes several
actions on the objects and learns what are the effects of those actions
in the environment conditioned to the shape features of the objects. A
probabilistic graphical model that encodes the joint probability
distribution of object features, actions and effects can be learned from
data using structure learning techniques. This model can then be used to
support several cognitive skills: action planning, action prediction,
object selection, action recognition and imitation.

\textbf{Related Papers:}

\begin{itemize}
\item
  L. Jamone, E. Ugur, A. Cangelosi, L. Fadiga, \ul{A. Bernardino}, J.
  Piater and J. Santos-Victor (2018), ``Affordances in Psychology,
  Neuroscience and Robotics: A Survey'', \emph{IEEE Transactions on
  Cognitive and Developmental Systems,} 10(1), pp 4-25.
  \url{https://doi.org/10.1109/TCDS.2016.2594134}
\item
  L. Montesano, M. Lopes, \ul{A. Bernardino}, J. Santos-Victor (2007),
  ``Modeling Affordances using Bayesian networks'', \emph{IEEE/RSJ
  International Conference on Intelligent Robots and Systems (IROS
  2007),} pp 4102-4107. \url{https://doi.org/10.1109/IROS.2007.4399511}
\end{itemize}

\clearpage

\textbf{\textsc{\ul{Learning Language Associations of Robotic Actions}}}

\textbf{Authors:} G. Salvi, L. Montesano, \ul{A. Bernardino}, J.
Santos-Victor

\textbf{Year:} 2012

\textbf{Title:} ``Language Bootstrapping: Learning Word Meanings From
Perception--Action Association''

\textbf{Journal:} \emph{IEEE Transactions on Systems, Man, and
Cybernetics, Part B: Cybernetics} 42(3), 660-671.

\textbf{Domain:} Bioinspired and Learning Robots, Human-Robot
Interaction

\textbf{IF:} 6.2

\textbf{Citations:} 22/63

\textbf{DOI:} \url{https://doi.org/10.1109/TSMCB.2011.2172420}

\textbf{Description:} Before the emergence of Large Language Models
(e.g. ChatGPT) and Visual-Language Learning (e.g. image captioning),
this paper has shown empirically that a robot agent can learn to
associate words to the entities of a robotic task: objects, actions and
effects. Learning is obtained by exploiting the co-occurrences between
the spoken words of a narration of the task, and the simultaneously
observed characteristics of the task measured by other perceptual
modalities. In combination with object's affordance models (see previous
work), the robot can learn the concepts of verbs associated to the robot
actions, names for objects, and adjectives to characteristics of the
objects and the actions. The model is learned in the form of a Bayesian
Network to represent the inherent the probabilistic nature of the
problem. Once learned, the model can be used to instruct the robot to
perform actions related to the objects present in the environment.
Through probabilistic reasoning and contextual cues the robot can
resolve ambiguous cases and incomplete instructions. This model had a
key role in the cognitive architecture of the Poeticon++ project, in
which the iCub robot interacted with humans in natural language to
execute manipulation tasks.

\textbf{Related Papers:}

\begin{itemize}
\item
  V. Krunic, G. Salvi, \ul{A. Bernardino}, L. Montesano, J.
  Santos-Victor (2009), ``Affordance Based Word-To-Meaning
  Association'', \emph{IEEE International Conference on Robotics and
  Automation (ICRA 2009)}, pp 806--811.
  \url{https://dl.acm.org/doi/10.5555/1703435.1703566}
\item
  Antunes, L. Jamone, G. Saponaro, \ul{A. Bernardino}, R. Ventura
  (2016), ``From Human Instructions to Robot Actions: Formulation of
  Goals, Affordances and Probabilistic Planning'', \emph{IEEE
  International Conference on Robotics and Automation (ICRA 2016),} pp
  5449-5454. \url{http://doi.org/10.1109/ICRA.2016.7487757}
\end{itemize}

\clearpage

\textbf{\textsc{\ul{Foveal Vision for Robots}}}

\textbf{Authors:} V. Javier Traver, \ul{A. Bernardino}

\textbf{Year:} 2010

\textbf{Title:} ``A review of log-polar imaging for visual perception in
robotics: A survey''

\textbf{Journal:} \emph{Robotics and Autonomous Systems,} 58(4),
378-398.

\textbf{Domain:} Bioinspired and Learning Robots, Vision for Decision
and Control

\textbf{IF:} 4.6

\textbf{Citations:} 134/210

\textbf{DOI:} \url{https://doi.org/10.1016/j.robot.2009.10.002}

\textbf{Description:} Foveal Vision mimics the distribution of
photoreceptors in the human retina and is characterized by a high
density in the center and a decaying resolution towards the periphery.
Coupled with eye movements, this representation allows to have
simultaneously a high-resolution area and a wide field of view with a
limited number of pixels, which are important characteristics of visual
systems for robots with limited computational resources.

The log-polar mapping is a popular mathematical representation of this
type of images, having interesting computational properties such as
equivariance to centered rotations and scaling. This paper reviews the
main representations, computational methods, and properties of log-polar
images and describes a large set of works that applied log-polar image
processing in robot vision, highlighting their advantages in computation
time and accuracy with respect to the conventional images. The paper
presents in a unified way a set of original works from myself at Vislab
where I show the advantages of foveal vision in robot eye-head control,
particularly in tracking, segmentation, and depth perception of central
objects.

\textbf{Related Papers:}

\begin{itemize}
\item
  \ul{A. Bernardino}, J. Santos-Victor (1999), ``Binocular Tracking:
  Integrating Perception and Control'', \emph{IEEE Transactions on
  Robotics and Automation,} (15)6, 1080-1094.
  \url{https://doi.org/10.1109/70.817671}
\item
  \ul{A. Bernardino}, J. Santos-Victor (1998) ``Visual Behaviours for
  Binocular Tracking'', \emph{Robotics and Autonomous Systems,} (25)3-4,
  137-146. \url{https://doi.org/10.1016/S0921-8890(98)00043-8}
\end{itemize}

\clearpage

\textbf{\textsc{\ul{Convex Methods for Learning Object Detectors in
Images}}}

\textbf{Authors:} R. Cabral, F. de La Torre, J. Costeira, \ul{A.
Bernardino}

\textbf{Year:} 2015

\textbf{Title:} ``Matrix Completion for Weakly-supervised Multi-label
Image Classification''

\textbf{Journal:} \emph{IEEE Trans. on} \emph{Pattern Analysis and
Machine Intelligence,} 37(1), pp 121-135

\textbf{Domain:} Vision for Decision and Control

\textbf{IF:} 23.6

\textbf{Citations:} 149/429

\textbf{DOI:} \url{https://doi.org/10.1109/TPAMI.2014.2343234}

\textbf{Description:} The most common methodology to learn object
detectors from image data is through supervised learning. This
methodology requires a dataset with annotated bounding boxes around the
objects to detect, which is typically done by laborious manual work on
large image sets. To reduce the amount of manual work, many authors have
proposed methods in which only the presence or absence of objects in an
image is provided (weak labels). Then, the method must learn from the
co-occurrence of object features in many images with coherent labels.
However, existing methods used combinatorial approaches that make the
learning process very slow. Instead, we proposed a convex method, that
exploits the feature-based histogram representation of images as the
linear combination of the histograms of the composing objects. With this
formulation, the problem can be solved efficiently through
matrix-completion with rank regularization, that is robust to noise in
the images and in the labeling process, as well as to moderate amounts
of missing labels and data. Experimental validation on several data sets
shows that our method outperforms state-of-the-art classification
algorithms, while effectively capturing each class appearance.

\textbf{Related Papers:}

\begin{itemize}
\item
  R. Cabral, F. De la Torre, J. Costeira, \ul{A. Bernardino} (2013),
  ``Unifying Nuclear Norm and Bilinear Factorization Approaches for
  Low-rank Matrix Decomposition'', \emph{IEEE International Conference
  on Computer Vision (ICCV 2013),} pp 2488-2495.
  \url{https://doi.org/10.1109/ICCV.2013.309}
\item
  R. Cabral, F. De la Torre, J. Costeira, \ul{A. Bernardino} (2011),
  ``Matrix Completion for Multi-label Image Classification'',
  \emph{International Conference on Neural Information Processing
  Systems} \emph{(NIPS 2011)}, pp 190-198.
  \url{https://dl.acm.org/doi/10.5555/2986459.2986481}
\end{itemize}

\clearpage

\textbf{\textsc{\ul{Computer Vision in Transportation Systems and
Autonomous Vehicles}}}

\textbf{Authors:} J. Melo, A. Naftel, \ul{A. Bernardino}, J.
Santos-Victor

\textbf{Year:} 2006

\textbf{Title:} ``Detection and Classification of Highway Lanes Using
Vehicle Motion Trajectories''

\textbf{Journal:} \emph{IEEE Transactions on Intelligent Transportation
Systems,} 7(2), 188-200.

\textbf{Domain:} Vision for Decision and Control

\textbf{IF:} 8.5

\textbf{Citations:} 112/184

\textbf{DOI:} \url{https://doi.org/10.1109/TITS.2006.874706}

\textbf{Description:} The automation of vehicles and road infrastructure
has witnessed a large penetration of computer vision and machine
learning methods in the last decade. We have applied some methods from
robot vision in applications to transportation systems, in particular
the detection and tracking of vehicles and persons. In the first paper
of this line of work, we have used low-level tracking methods to compute
vehicle trajectories, that, accompanied by clustering and
indexing/retrieval methods, were shown to: automatically identify the
lanes of traffic, compute statistics of traffic in highways, and detect
anomalous trajectories. Experimental results are presented to show the
real-time application of this approach to multiple views obtained by an
uncalibrated pan-tilt-zoom traffic camera monitoring the junction of two
busy intersecting highways. In follow-up papers, we have worked on ways
to improve object detection in autonomous driving applications.

\textbf{Related Papers:}

\begin{itemize}
\item
  D. Ribeiro, J. Nascimento, \ul{A. Bernardino}, G. Carneiro (2017),
  ``Improving the performance of pedestrian detectors using
  convolutional learning'', \emph{Pattern Recognition,} 61, pp 641--649.
  \url{https://doi.org/10.1016/j.patcog.2016.05.027}
\item
  M. Taiana, J. Nascimento, \ul{A. Bernardino} (2015), ``On the purity
  of training and testing data for learning: The case of pedestrian
  detection'', \emph{Neurocomputing}, 150A, pp 214-226.
  \url{https://doi.org/10.1016/j.neucom.2014.09.055}
\end{itemize}

\clearpage

\textbf{\textsc{\ul{Airborne Image Analysis for Maritime Surveillance}}}

\textbf{Authors:} R. Ribeiro, G. Cruz, J. Matos, \ul{A. Bernardino}

\textbf{Year:} 2019

\textbf{Title:} ``A Data Set for Airborne Maritime Surveillance
Environments''

\textbf{Journal:} \emph{IEEE Trans. on Circuits and Syst. for Video
Technology,} 29(9), pp 2720-2732

\textbf{Domain:} Vision for Decision and Control

\textbf{IF:} 8.4

\textbf{Citations:} 82/110

\textbf{DOI:} \url{https://doi.org/10.1109/TCSVT.2017.2775524}~~

\textbf{Description:} Maritime surveillance is a key activity for many
countries. It is important to assure the safe and secure use of the
oceans for transportation and trade. It allows the control of fisheries
to guarantee the protection of resources and ecosystems. Maritime
surveillance also ensures that environmental regulations are applied,
preventing oil spill and bilge dumping that have a severe impact on
fauna, flora and also coast human populations. In the last decade,
Unmanned Aerial Vehicles (UAVs) have seen a huge increase not only in
its deployment but also in its capabilities. Right now, UAVs offer
promising technologies to aid remote sensing and oceanic surveillance.
While traditional aircraft are equipped with heavy radars, UAVs normally
have only light weight passive electro-optical sensors. This paper
presents a data set with surveillance imagery over the sea captured by a
small size UAV. Additionally, using standard evaluation frameworks, we
establish a baseline of results for automated image and video analysis
methods, using algorithms developed by the authors, which are better
adapted to the maritime environment.

\textbf{Related Papers:}

\begin{itemize}
\item
  G. Cruz, \ul{A. Bernardino} (2019), ``Learning Temporal Features for
  Detection on Maritime Airborne Video Sequences using Convolutional
  LSTM'', \emph{IEEE Transactions on Geoscience and Remote Sensing,}
  57(9), pp 6565-6576.
  \url{https://doi.org/10.1109/TGRS.2019.2907277}~~~
\item
  G. Cruz, \ul{A. Bernardino} (2016), ``Aerial detection in maritime
  scenarios using convolutional neural networks'', \emph{International
  Conference on Advanced Concepts for Intelligent Vision Systems (ACIVS
  2016),} Lecture Notes in Computer Science, Springer, vol 10016, pp
  373-384. \url{https://doi.org/10.1007/978-3-319-48680-2_33}
\end{itemize}

\clearpage

\textbf{\textsc{\ul{Human Recognition and Tracking}}}

\textbf{Authors:} A. Nambiar, \ul{A. Bernardino}, J. Nascimento

\textbf{Year:} 2019

\textbf{Title:} ``Gait-Based Person Re-identification: A Survey'',

\textbf{Journal:} \emph{ACM Computing Surveys (CSUR)} 52(2), 33, pp
1-34.

\textbf{Domain:} Vision for Decision and Control, Human-Robot
Interaction

\textbf{IF:} 16.6

\textbf{Citations:} 90/146

\textbf{DOI:} \url{https://doi.org/10.1145/3243043}~

\textbf{Description:} Person Re-Identification is the ability to
recognize the same person at two different instances of time. It has
important applications in surveillance systems and human-robot
interaction. In our work we have developed several methodologies for
person re-identification with computer vision using appearance, shape,
facial, and gait features. In this paper we perform and extensive survey
of gait-based person reidentification. Gait features are attractive
because they are persistent across extended periods of time and are hard
to fake, allowing long-term re-identification. However, the acquisition
of good-quality measures of a person's motion patterns in unconstrained
environments has proved very challenging in practice. Existing
technology (video cameras) suffers from changes in viewpoint, daylight,
clothing, accessories, and other variations in the person's appearance.
We identify several relevant dimensions of the problem and provide a
taxonomic analysis of the current state of the art. We present some of
our own work in this field, discussing the levels of performance
achievable with the current technology. Finally, we give a perspective
of the most challenging and promising directions of research for the
future.

\textbf{Related Papers:}

\begin{itemize}
\item
  D. Figueira, L. Bazzani, H. Q. Minh, M. Cristani, \ul{A. Bernardino},
  V. Murino (2013), ``Semi-supervised Multi-feature Learning for Person
  Re-identification'', \emph{IEEE International Conference on Advanced
  Video and Signal Based Surveillance,} pp 111-116.
  \url{https://doi.org/10.1109/AVSS.2013.6636625}
\item
  Nambiar, \ul{A. Bernardino}, J. Nascimento (2015), ``Shape Context for
  soft biometrics in person re-identification and database retrieval'',
  \emph{Pattern Recognition Letters,} 68(2), pp 297--305.
  \url{https://doi.org/10.1016/j.patrec.2015.07.001}
\end{itemize}

\clearpage

\textbf{\textsc{\ul{3D Model-Based Object Tracking}}}

\textbf{Authors:} N. Pessanha-Santos, V. Lobo, \ul{A. Bernardino}

\textbf{Year:} 2020

\textbf{Title:} ``Two-stage 3D model-based UAV pose estimation: A
comparison of methods for optimization''

\textbf{Journal:} \emph{Journal of Field Robotics,} 37(4), pp 580-605.

\textbf{Domain:} Vision for Decision and Control

\textbf{IF:} 8.3

\textbf{Citations}: 12/22

\textbf{DOI:} \url{https://doi.org/10.1002/rob.21933}~

\textbf{Description:} 3D Model-Based Object Tracking is a methodology
that exploits available 3D models of objects to detect and track their
6D pose (position and orientation) with respect to a camera sensor. It
has important applications in robot navigation and control, and object
grasping and manipulation. There are two main approaches to the problem:
inverse and direct methods. Inverse methods try to solve the problem of
transforming the 2D measurements obtained in the image to the 6D pose of
the object. Direct methods, instead, compute the likelihood of the 6D
object pose given the 2D data. Direct methods tend to be more robust to
moderate amounts of occlusion and variability in the scene
(illumination, clutter, background), but face two main problems: (i) the
true likelihood function is unknown, so approximations must be carefully
designed for the problem at hand, and (ii) even approximate likelihood
functions are highly irregular, so sampling methods are required for its
optimization which brings computational complexity issues. In this
paper, we propose and compare several likelihood functions and
optimization methods to the problem of 3D model-based tracking. The
framework is applied to the outdoor pose estimation of a fixed-wing UAV
for autonomous landing in a Fast Patrol Boat (FPB).

\textbf{Related Papers:}

\begin{itemize}
\item
  N. Pessanha-Santos, V. Lobo, \ul{A. Bernardino} (2020), ``Directional
  statistics for 3D Model-Based UAV Tracking'', \emph{IEEE Access,} 8,
  pp 33884-33897. \url{https://www.doi.org/10.1109/ACCESS.2020.2973970}~
\item
  N. Pessanha-Santos, V. Lobo and \ul{A. Bernardino} (2019), ``Unmanned
  Aerial Vehicle Tracking Using a Particle Filter Based Approach'',
  \emph{IEEE Underwater Technology (UT 2019),} pp 1-10.
  \url{https://doi.org/10.1109/UT.2019.8734465}
\end{itemize}

\clearpage

\textbf{\textsc{\ul{Robots to Study Biological Systems}}}

\textbf{Authors:} A. John, C. Aleluia, A. J. Van Opstal, A. Bernardino

\textbf{Year:} 2021

\textbf{Title:} ``Modelling 3D saccade generation by feedforward optimal
control''

\textbf{Journal:} \emph{PLoS Computational Biology} 17(5), e1008975, pp
1-35.

\textbf{Domain:} Bioinspired and Learning Robots

\textbf{IF:} 4.3

\textbf{Citations:} 11/20

\textbf{DOI:} \url{https://doi.org/10.1371/journal.pcbi.1008975}

\textbf{Description:} Using robotic models to test theories on the
behavior of humans or animals can help understand many aspects of
intelligence and control in natural systems. This work resulted from a
collaboration between a team of roboticist and a team of
neuroscientists, both studying the human oculomotor system from
different perspective. The paper focus on how model-based optimal
control principles can explain stereotyped human oculomotor behaviors
(Listing's Law, Saccade dynamics), through simulations in a realistic
model of the human eye with a cable-driven actuation system. ~We
verified the contribution of different terms in the cost optimization
functional to realistic 3D saccade behavior, and identified four
essential terms: total energy expenditure by the motors, movement
duration, gaze accuracy, and the total static force exerted by the
muscles during fixation. Our findings suggest that Listing's law, as
well as the saccade dynamics and their trajectories, may all emerge from
the same common mechanism that aims to optimize speed-accuracy trade-off
for saccades, while minimizing the total muscle force during eccentric
fixation. Other papers in this line of research have focused on learning
aspects and studied the visual systems of other animal species.

\textbf{Related Papers:}

\begin{itemize}
\item
  R. J. Alitappeh, A. John, B. Dias, A. J. van Opstal and \ul{A.
  Bernardino} (2024) "Emergence of human oculomotor behavior in a
  cable-driven biomimetic robotic eye using optimal control,"
  in~\emph{IEEE Transactions on Cognitive and Developmental Systems},
  6(4):1546-1560. \url{https://doi.org/10.1109/TCDS.2024.3376072}
\item
  X. Liu, M. D. Loring, L. Zunino, K. E. Fouke, F. A. Longchamp, \ul{A.
  Bernardino}, A. J. Ijspeert, E. A. Naumann (2025) ``Artificial
  Embodied Circuits Uncover Neural Architectures of Vertebrate
  Visuomotor Behaviors'', \emph{Science Robotics,} 10,
  \url{https://doi.org/10.1126/scirobotics.adv4408}
\end{itemize}

\clearpage

\subsubsection{International
Collaborations}\label{international-collaborations}

Among my 460+ co-authors in scientific publications, most of them are
international. In this section I list 10 of my closer and more
established international collaborations.

\begin{enumerate}
\def\labelenumi{\arabic{enumi}.}
\item
  \textbf{Lorenzo Jamone, 29 co-authored works}, \textbf{3 joint PhD
  students}, Associate Professor in Robotics \& AI at the Department of
  Computer Science of University College London (UCL), \textbf{UK}.
\item
  \textbf{Asim Smailagic}, \textbf{15 co-authored works}, \textbf{3
  joint PhD students}, Research Professor at CMU, Director of Laboratory
  for Interactive Computer Systems and Wearable Computers, \textbf{US}.
\item
  \textbf{Daniel Siewiorek}, \textbf{14 co-authored works}, \textbf{2
  joint PhD students}, Buhl University Professor at CMU, Human-Computer
  Interaction Institute, \textbf{US}.
\item
  \textbf{Auke Jan Ijspeert}, \textbf{11 co-authored works}, \textbf{2
  joint PhD} \textbf{students}, full professor at EPFL (the Swiss
  Federal Institute of Technology at Lausanne), head of the Biorobotics
  Laboratory (BioRob), \textbf{Switzerland}.
\item
  \textbf{Cecilia Laschi}, \textbf{10 co-authored works}, full professor
  at the National University of Singapore, College of Design and
  Engineering, \textbf{Singapore}.
\item
  \textbf{A. J. Van Opstal}, \textbf{9 co-authored works}, \textbf{1
  joint PhD student}, professor at Donders Centre for Neuroscience and
  Donders Institute for Brain, Cognition and Behaviour,
  \textbf{Netherlands}.
\item
  \textbf{Paolo Dario}, \textbf{9 co-authored works}, professor at
  Scuola Superiore Sant'Anna (SSSA), Pisa, and Founding Coordinator of
  the Center of Micro-BioRobotics@SSSA of the Italian Institute of
  Technology (IIT) \textbf{Italy}.
\item
  \textbf{Fernando de La Torre}, \textbf{8 co-authored works}, \textbf{2
  joint PhD students}, Research Associate Professor at CMU, Robotics
  Institute, \textbf{US}.
\item
  \textbf{Giampiero Salvi}, \textbf{6 co-authored works}, \textbf{1
  joint PhD student}, professor at the Norwegian University of Science
  and Technology, \textbf{Norway}, and Associate Professor with KTH
  Royal Institute of Technology, \textbf{Sweden}.
\item
  \textbf{Jodi Forlizzi, 5 co-authored works, 1 joint PhD student,}
  Herbert A. Simon Professor in Computer Science and~HCII at CMU,
  Human-Computer Interaction Institute, \textbf{US}.
\end{enumerate}

\subsection{Scientific Project Grants}\label{scientific-project-grants}

I have participated in more than \textbf{40 scientific research
projects} (national and international) involving a budget of more than
\textbf{7.5M€ at IST}. While participating as \textbf{global
coordinator} (8 projects) or \textbf{coordinator at the participating
institution} (8 projects), I have secured a budget of more than
\textbf{2.5M€ for IST}. A complete description of my projects is
provided in Appendix \ref{project-grants}.

I highlight my participations in the following projects:

\begin{itemize}
\item
  I was the principal investigator of the CMU-PT project \textbf{AHA}
  (2014-2018) where I initiated strong academic links with four of my
  top international collaborators: Prof. Asim Smailagic, Prof. Dan
  Siewiorek, Prof. Fernando de la Torre, and Prof. Jodi Forlizzi. In
  this project we developed a novel robotic assistance platform
  (\textbf{PEPE}) designed to support healthy lifestyle, sustain active
  aging, and support those with motor deficits. \textbf{PEPE} has a
  \textbf{national and a European patent}. The platform is being used in
  a dozen of institutions in Portugal and has been exploited more
  recently to develop projects related to cognitive stimulation of
  senior adults.
\item
  I was the principal investigator of the national project
  \textbf{HAVATAR} (2021-2024), that deepened collaborations with an
  important partner in the health sector, \textbf{Hospital da Luz},
  GLSMED Learning Health, S.A. and studied the application of
  technologies in the health sector, namely the use of
  \textbf{teleoperation robots} to allow remote visits in hospitals,
  motivated by the pandemics context. It also continued the development
  of virtual and robotic coaches initiated in the AHA projects for
  \textbf{stroke rehabilitation}.
\item
  I was the principal investigator of the national projects
  \textbf{FIREFRONT} (2019-2023) and \textbf{VOAMAIS} (2019-2023), and
  local PI of the \textbf{SEAGULL} (2013-2015) and \textbf{ELEVAR}
  (2017-2019) projects, that developed \textbf{airborne computer vision}
  and robotics technology (deployed in drones) for the
  \textbf{monitoring of critical infrastructures} (bridges, buildings,
  dams), \textbf{wildfires in the national forest}, and for the
  \textbf{surveillance of maritime events} like illegal naval traffic
  and pollution spots. These projects developed strong links with the
  \textbf{Naval School} and the \textbf{Air Force Academy}, as well as
  with national technological companies like \textbf{Tekever},
  \textbf{UAVision}, and \textbf{Critical Software}.
\item
  I was the local PI of the EU project \textbf{ORIENT} (2019-2022) in
  which ISR-IST participated as a beneficiary of the \textbf{ERC Grant}
  of the neuroscientist Prof. John Van Opstal, \emph{Stichting Radboud
  Universiteit, Nijmegen} (NL). In this project Prof. John Van Opstal
  studied the Head-Eye orienting responses of animals in rapid object
  identification. Given the current limitations in the use of animal
  studies in Europe, Prof. John Van Opstal invited ISR-IST to develop
  robotic analogues of a human eye to test several theories on rapid eye
  movement control (saccades). We developed the \textbf{first cable
  driven mechanical eye with 6 degrees of freedom}, mimicking the main
  biomechanical aspects of the human eye, that was used usefully to
  \textbf{test optimal control and reinforcement learning theories} in
  the emergence of saccade eye movements.
\item
  I was the Technical Manager of the EU project \textbf{HANDLE}
  (2009-2013) closely supporting the PI in daily technical decisions on
  the project. We studied how humans perform the manipulation of objects
  and used this knowledge to replicate grasping and skilled
  \textbf{in-hand movements with an anthropomorphic artificial hand}. We
  used the most advanced robot hand in the world, the
  \emph{\textbf{Shadow Robot Hand}}, and produced pioneering works in
  dexterous manipulation using hand synergies and objects' affordances.
  In this project the PI and I jointly supervised a PhD student in a
  co-tutelle agreement between IST and UPMC, Paris.
\item
  I was one of the key team members in the EU Project \textbf{ROBOTCUB}
  (2004-2010) that resulted in the construction of the humanoid robot
  iCub. The iCub is the most complex humanoid robot worldwide (in terms
  of degrees of freedom) and is an open platform to study developmental
  robotics, i.e. how humans develop \textbf{advanced sensorimotor
  coordination, learning, perception, interaction, planning, and
  execution skills}. Vislab has the only copy of the iCub in Portugal,
  as part of the RBCog lab, which has been used to support the
  graduation of many PhD students. It created a community of users and
  spin-off projects that paved the way for the growth of research in the
  \textbf{humanoid robotics} field.
\item
  I was the responsible for the cross-platform computational
  architecture of the EU project \textbf{NARVAL} (1996-2001). I
  highlight this project as it was my first important contact with a
  research and development project while I was a PhD student. The
  project was globally coordinated by Prof. Santos-Victor of Vislab
  ISR/IST and developed \textbf{computer vision technologies for the
  control of underwater autonomous vehicles}. I developed
  \textbf{algorithms for tracking planar surfaces underwater} for
  stabilization and docking of AUVs. Also, I developed a
  \textbf{distributed software architecture} for the integration of
  several computational modules and interoperability with the software
  of other partners.
\end{itemize}

\subsection{Recognition and
Visibility}\label{recognition-and-visibility}

My scientific activities have been very well appreciated in the research
community. I had \textbf{7 distinctions related to publications} (best
paper, best poster, honourable mentions, and alike), and \textbf{6
distinctions related to my students' theses} (best thesis, outstanding
thesis). I was a \textbf{finalist} (Top 5 among 85 participants) of the
\textbf{Smart Ageing Prize}, AAL Programme, AAL Europe -- ICT for Ageing
Well, 2018. I am a \textbf{Senior Member of IEEE Robotics and Automation
Society} and member of \textbf{ELLIS Society}, associated with the
Lisbon Unit (\textbf{LUMLIS}). I am regularly invited to seminars,
keynote talks, and plenary session in international and national events
to present the results of my research, totalling more than \textbf{40
events}. Complete lists of my awards, invited talks, and society
memberships is provided in Appendix \ref{recognition-and-visibility-1}.
